\section{Discussion}\label{sec:discussion}

Thresholds selected on point estimates or percentile-bootstrap bounds often claimed more evidence than held. These are
call-averaged ratios, not the local ratios of ClinGen calibrations, so this does not show that published PP3/BP4
thresholds are wrong, only that point estimates and bootstrap bounds can overstate evidence at single-gene sample sizes.
Confidence bounds for the ratio, exact or Katz, held in at least 97.9\% of draws; the exact certificate is the
conservative choice that keeps a guarantee at every sample size. \textsc{CertiVar} applies to genes with a functional
assay: to audit whether a predictor's claimed evidence strength holds against the assay, and to call the variants the
assay missed at assayed residues, for predictors whose scores pass the mixture check. The cost is conservatism:
certified evidence was mostly supporting-level, and PAC thresholds abstained on 58\% to 82\% of variants.

\paragraph{Limitations.}
The labels are assay classes, not clinical pathogenicity \cite{Brnich2020}, and the ACMG/AMP levels are reused as a
reporting scale; clinical use requires an assay validated against clinical classifications, as its authors did for BRCA1
\cite{Findlay2018} and BRCA2 \cite{Huang2025}. Our BRCA2 labels come from the assay's controls (\cref{tab:app_labels}).
A certified level describes the calls as a class and weights residues equally, so a borderline call can carry less
evidence than its class. The guarantees do not cover intermediate, splice-region or out-of-region variants;
intermediate variants are a hazard, because whether a VUS would score intermediate is unknown before it is assayed, and
19.6\% of them were called LoF. The mixture check tests only a necessary condition and missed most VUS shifts. Results
cover five cancer-predisposition genes. No predictor was trained on these assay labels, but REVEL was trained on
labelled variants and the AlphaMissense classes were set on ClinVar, so both may overlap our evaluation sets
\cite{Ioannidis2016,Grimm2015}; ESM-2, trained without variant labels, gave the same coverage and validity.

\paragraph{Conclusion.}\label{sec:conclusion}
\textsc{CertiVar} turns a variant-effect predictor and a few hundred assayed residues of a gene into LoF, functional or
no calls with finite-sample guarantees on per-class error rates and on call-averaged likelihood ratios for the assay
after threshold selection, without an estimate of how common damaging variants are, and checks whether a target set can
inherit them. Code, data and results (release v1.1.1):\newline\url{https://github.com/aaryavvatts-lab/certivar}.
